\section{Practica 1: Cuatro Leds}

\subsection{Introducción}
\noindent	
	Esta práctica tratará de encender cuatro focos LED, usando una PIC16F628A, una protoboard, cuatro resistencias y  cuatro focos LED. Es una práctica introductoria que tiene como finalidad presentar el uso más básico y elemental del PIC16F628A. Un PIC es un circuito integrado programable (Programmable Integrated Circuit), el cual contiene todos los componentes para poder realizar y controlar una tarea, por lo que se denomina como un microcontrolador.

\subsection{Materiales}
	\begin{itemize}
		\item Microcontrolador PIC16F628A.
		\item 1 Protoboard
		\item Cristal de cuarzo 4MHz (Oscilador XT).
		\item 4 Resistencias de 220 $\Omega$ (para los LEDs).
		\item 4 LEDs individuales.
		\item Fuente de alimentación de 5V DC
		\item 2 capacitores 22pF
		\item Cables
	\end{itemize}

\subsection{Software}
	\begin{itemize}
		\item Proteus 8.0
		\item MPLAB X IDE
		\item PICkit 3 Programmer App and Scripting Tool v3.10
	\end{itemize}

\subsection{Circuito}
	\begin{figure}[H]
		\centering
		\includegraphics[width = 0.8 \textwidth]{PROY_01.png}
		\caption{Circuito exportado de Proteus}
		\label{fig:proy01}
	\end{figure}

\subsection{Desarrollo}
\noindent
	Empezamos esta práctica haciendo el diagrama en proteus (Figura \ref{fig:proy01}), para posteriorment	e, en el software de MPLAB, en código C, hacer el programa .hex, el cual probamos en una 	simulación en proteus.\\
Una vez hecho esto, procedimos a conectar en nuestro Protoboard los materiales necesario	s.

	\begin{figure}[H]
		\centering
		\includegraphics[width = 0.8 \textwidth]{.png}
		\caption{Conexion de componentes}
		\label{fig:concomponentes}
	\end{figure}
\noindent	
	Programamos el PIC16F628A con ayuda de nuestro PicKit 3, y lo conectamos a nuestro protoboard.

	%aqui va otra figura
\noindent	
	Este fue nuestro primer intento, y no prendió, nos dimos cuenta de que estaba mal conectado en el puenteo.

	%otra figura
\noindent	
	Ya al corregir todo volvimos a probar, y ya nos salió.

	%otra figura

\subsection{Codigo en C}
	\begin{verbatim}
		#pragma config FOSC = XT        // Oscillator Selection bits 
		(XT oscillator: Crystal/resonator on RA6/OSC2/CLKOUT and RA7/OSC1/CLKIN)
#pragma config WDTE = OFF       // Watchdog Timer Enable bit (WDT disabled)
#pragma config PWRTE = ON       // Power-up Timer Enable bit (PWRT enabled)
#pragma config MCLRE = ON       // RA5/MCLR/VPP Pin Function Select bit 
(RA5/MCLR/VPP pin function is MCLR)
#pragma config BOREN = ON       // Brown-out Detect Enable bit (BOD enabled)
#pragma config LVP = OFF        // Low-Voltage Programming Enable bit 
(RB4/PGM pin has digital I/O function, HV on MCLR must be used for programming)
#pragma config CPD = OFF        // Data EE Memory Code Protection bit 
(Data memory code protection off)
#pragma config CP = OFF         // Flash Program Memory Code Protection bit 
(Code protection off)

// #pragma config statements should precede project file includes.
// Use project enums instead of #define for ON and OFF.

#include <xc.h>
void main(void) {
    TRISA = 0b11111;
    TRISB = 0b10101010;
    
    PORTA = 0x0;
    PORTB = 0x0;
    
    while(1){
        PORTB = 0b1010101;
    }
    return;
}

	\end{verbatim}

\subsection{Circuito cuando se ejecuta la simulacion}
	\begin{figure}[H]
		\centering
		\includegraphics[width = 0.8 \textwidth]{PROY_01_sim.png}
		\caption{Circuito exportado de Proteus}
		\label{fig:proy01sim}
	\end{figure}

\subsection{Conclusion}
\noindent	
	En esta práctica se logró implementar correctamente el encendido de cuatro LEDs mediante el uso de un microcontrolador PIC16F628A, integrando tanto la programación en lenguaje C como el armado del circuito en protoboard. A través de este ejercicio se comprendió la importancia de la configuración de los puertos de entrada y salida, así como el control digital de señales para manipular dispositivos electrónicos básicos.

\noindent	
Además, se reforzaron habilidades en la interpretación de diagramas eléctricos y en la correcta conexión de componentes, evitando errores comunes como cortocircuitos o conexiones incorrectas. La programación permitió establecer una lógica de control sobre los LEDs, demostrando cómo el software interactúa directamente con el hardware.

